\chapter{Experimental Data and Processing}
\label{app:data}

\section{Measurements}

The electrical tables give instrument readings for completed fixed-output and extended-load tests. Blank fields mark missing measurements.

Seven V2 secondary-output load captures provide AC-supply and electronic-load readings used to calculate secondary power and estimate efficiency as described in \cref{sec:method-efficiency}. The associated thermal-region readings are described in \cref{sec:method-thermal}.

\section{Calculations and Graph Generation}

Measured values were kept in the electrical tables. The equations in
Appendix~\ref{app:equations} were used to calculate output power,
average-voltage error, initial efficiency estimates, and the input-output power
difference. Efficiency requires a compatible DC average, load current,
and logged input power. Incomplete rows, unmatched input-power readings,
and rows flagged for power disagreement are excluded. Some selected rows
still have differing instrument readings, which limit the accuracy of the
efficiency estimates. PF and THD readings remain linked to their instruments;
they cannot be treated as one calibrated measurement across all tests.

AC-coupled detail captures and DC operating-point measurements are listed separately. The spread of \texttt{Vavg} statistics describes variation between average readings; ripple requires an instantaneous peak-to-peak voltage from a defined capture. The thermal results use the maximum surface temperature displayed by the camera. A junction-temperature estimate or continuous rating would also need component identification, a thermal model, and a stabilized run.

The raw waveform exports were plotted with their original amplitudes
and offsets. Uncertainty in channel calibration or the measured circuit
node is stated with the result. 

\section{Electrical Measurements and Calculated Results}
The tables list voltages, load, power-quality readings, firmware version,
and calculations for each completed test. DC
operating points, AC-detail captures, and load-removal tests are kept
separate because the quantities and acquisition conditions differ.
% Automatically rebuilt from original records by rebuild_results.py.
\begin{landscape}
\begingroup
\scriptsize
\setlength{\tabcolsep}{3pt}
\renewcommand{\arraystretch}{1.12}
\begin{longtable}{@{}L{0.2038\linewidth} L{0.0648\linewidth} L{0.0648\linewidth} L{0.0834\linewidth} L{0.0834\linewidth} L{0.0834\linewidth} L{0.0834\linewidth} L{0.0741\linewidth} L{0.0741\linewidth} L{0.0648\linewidth}@{}}
\caption{DC operating-point measurements at 240 V, 50 Hz, with calculated output power and voltage error. VBUS pp keeps the original workbook label.}\label{tab:appendix-dc-complete}\\
\toprule
Test ID & $V_\mathrm{set}$ (V) & $I_\mathrm{out}$ (A) & $V_\mathrm{BUS}$ (V) & Logged pp (V) & $V_\mathrm{sec}$ (V) & Sec. pp (V) & $P_\mathrm{in}$ (W) & $P_\mathrm{out}$ (W) & Error (\%)\\
\midrule
\endfirsthead
\toprule
Test ID & $V_\mathrm{set}$ (V) & $I_\mathrm{out}$ (A) & $V_\mathrm{BUS}$ (V) & Logged pp (V) & $V_\mathrm{sec}$ (V) & Sec. pp (V) & $P_\mathrm{in}$ (W) & $P_\mathrm{out}$ (W) & Error (\%)\\
\midrule
\endhead
\sourcefile{PDO5-NL-001} & 5 & 0.00 & 5.00727 & 1.396 & 33.0362 & 21.7449 & -- & 0.000 & 0.145\\
\sourcefile{PDO5-5W-001} & 5 & 1.00 & 5.02913 & 2.634 & 32.1636 & 27.2306 & 9.195 & 5.029 & 0.583\\
\sourcefile{PDO5-10W-2A-001} & 5 & 2.00 & 4.99630 & 2.255 & 28.5689 & 23.4159 & 14.779 & 9.993 & -0.074\\
\sourcefile{PDO5-MAX3A-001} & 5 & 3.00 & 5.02118 & 2.252 & 28.4363 & 24.1054 & 21.585 & 15.064 & 0.424\\
\sourcefile{PDO9-NL-001} & 9 & 0.00 & 9.02697 & 1.517 & 32.2001 & 26.9399 & -- & 0.000 & 0.300\\
\sourcefile{PDO15-NL-001} & 15 & 0.00 & 15.05370 & 1.092 & 25.3911 & 34.6820 & -- & 0.000 & 0.358\\
\sourcefile{PDO20-5W-001} & 20 & 0.25 & 20.04910 & 12.656 & 31.5934 & 12.6230 & 8.780 & 5.012 & 0.245\\
\sourcefile{PDO20-75W-001} & 20 & 3.70 & 19.99090 & 0.416 & 28.7198 & 17.0668 & 95.875 & 73.966 & -0.045\\
\sourcefile{PDO20-76W-38A-001} & 20 & 3.80 & 19.98130 & 0.452 & 28.7478 & 18.2891 & 98.990 & 75.929 & -0.093\\
\sourcefile{PDO20-80W-001} & 20 & 3.90 & 19.98320 & 0.453 & 28.7878 & 18.5766 & 95.875 & 77.934 & -0.084\\
\sourcefile{PDO20-80W-40A-001} & 20 & 4.00 & 19.97320 & 0.421 & 28.7255 & 17.1853 & 104.418 & 79.893 & -0.134\\
\sourcefile{PDO9-9W-1A-001} & 9 & 1.00 & 9.02079 & 1.291 & 28.1530 & 23.9798 & 10.795 & 9.021 & 0.231\\
\sourcefile{PDO9-18W-2A-001} & 9 & 2.00 & 9.01365 & 1.318 & 28.4514 & 23.9060 & 20.927 & 18.027 & 0.152\\
\sourcefile{PDO9-27W-3A-001} & 9 & 3.00 & 9.00057 & 1.220 & 27.7108 & 22.7355 & 31.101 & 27.002 & 0.006\\
\sourcefile{PDO15-15W-1A-001} & 15 & 1.00 & 15.01040 & 0.792 & 27.8648 & 24.9056 & 17.372 & 15.010 & 0.069\\
\sourcefile{PDO15-30W-2A-001} & 15 & 2.00 & 14.99490 & 0.767 & 28.5607 & 24.1054 & 34.071 & 29.990 & -0.034\\
\sourcefile{PDO15-45W-3A-001} & 15 & 3.00 & 14.98350 & 0.349 & 28.4765 & 20.3051 & 50.622 & 44.950 & -0.110\\
\bottomrule
\end{longtable}
\endgroup
\end{landscape}

% Automatically rebuilt from original records by rebuild_results.py.
\begingroup
\scriptsize
\setlength{\tabcolsep}{3pt}
\renewcommand{\arraystretch}{1.12}
\begin{longtable}{@{}L{0.2872\linewidth} L{0.0741\linewidth} L{0.0741\linewidth} L{0.1297\linewidth} L{0.1297\linewidth} L{0.1853\linewidth}@{}}
\caption{AC-coupled measurements. Several workbook values describe the spread of average-voltage readings, not instantaneous peak-to-peak ripple. CH2 and CH3 do not show absolute DC voltage.}\label{tab:ac-coupled-ripple-captures-complete}\\
\toprule
Test ID & Setting (V) & Load (A) & VBUS detail (mV) & Secondary detail (V) & Interpretation\\
\midrule
\endfirsthead
\toprule
Test ID & Setting (V) & Load (A) & VBUS detail (mV) & Secondary detail (V) & Interpretation\\
\midrule
\endhead
\sourcefile{PDO15-NL-AC-001} & 15 & 0.00 & 24.719 & 1.52727 & AC statistics/detail; no DC level\\
\sourcefile{PDO15-15W-1A-AC-001} & 15 & 1.00 & 27.380 & 1.52727 & AC statistics/detail; no DC level\\
\sourcefile{PDO15-30W-2A-AC-001} & 15 & 2.00 & 28.985 & 1.52727 & AC statistics/detail; no DC level\\
\sourcefile{PDO15-45W-3A-AC-001} & 15 & 3.00 & 28.985 & 1.52727 & AC statistics/detail; no DC level\\
\sourcefile{PDO5-5W-1A-AC-001} & 5 & 1.00 & 28.985 & 10.49900 & AC statistics/detail; no DC level\\
\sourcefile{PDO5-10W-2A-AC-001} & 5 & 2.00 & 28.985 & 10.49900 & AC statistics/detail; no DC level\\
\sourcefile{PDO5-15W-3A-AC-001} & 5 & 3.00 & 29.692 & 10.49900 & AC statistics/detail; no DC level\\
\sourcefile{PDO20-NL-AC-001} & 20 & 0.00 & 76.218 & 10.49900 & AC statistics/detail; no DC level\\
\sourcefile{PDO20-40W-2A-AC-001} & 20 & 2.00 & 76.218 & 10.49900 & AC statistics/detail; no DC level\\
\sourcefile{PDO20-60W-3A-AC-001} & 20 & 3.00 & 83.391 & 10.49900 & AC statistics/detail; no DC level\\
\sourcefile{PDO20-80W-4A-AC-001} & 20 & 4.00 & 83.391 & 10.49900 & AC statistics/detail; no DC level\\
\sourcefile{PDO20-94W-47A-AC-001} & 20 & 4.70 & 83.391 & 10.49900 & AC statistics/detail; no DC level\\
\bottomrule
\end{longtable}
\endgroup

% Automatically rebuilt from original records by rebuild_results.py.
\begingroup
\scriptsize
\setlength{\tabcolsep}{3pt}
\renewcommand{\arraystretch}{1.12}
\begin{longtable}{@{}L{0.3022\linewidth} L{0.1067\linewidth} L{0.1244\linewidth} L{0.1244\linewidth} L{0.0978\linewidth} L{0.1244\linewidth}@{}}
\caption{Initial efficiency estimates from eleven measurements.}\label{tab:efficiency-results-complete}\\
\toprule
Test ID & $I_\mathrm{out}$ (A) & $P_\mathrm{out}$ (W) & $P_\mathrm{in}$ (W) & $\eta$ (\%) & Difference (W)\\
\midrule
\endfirsthead
\toprule
Test ID & $I_\mathrm{out}$ (A) & $P_\mathrm{out}$ (W) & $P_\mathrm{in}$ (W) & $\eta$ (\%) & Difference (W)\\
\midrule
\endhead
\sourcefile{PDO5-5W-001} & 1.00 & 5.03 & 9.20 & 54.69 & 4.17\\
\sourcefile{PDO5-10W-2A-001} & 2.00 & 9.99 & 14.78 & 67.61 & 4.79\\
\sourcefile{PDO5-MAX3A-001} & 3.00 & 15.06 & 21.59 & 69.79 & 6.52\\
\sourcefile{PDO20-5W-001} & 0.25 & 5.01 & 8.78 & 57.09 & 3.77\\
\sourcefile{PDO20-75W-001} & 3.70 & 73.97 & 95.88 & 77.15 & 21.91\\
\sourcefile{PDO20-76W-38A-001} & 3.80 & 75.93 & 98.99 & 76.70 & 23.06\\
\sourcefile{PDO20-80W-40A-001} & 4.00 & 79.89 & 104.42 & 76.51 & 24.53\\
\sourcefile{PDO9-9W-1A-001} & 1.00 & 9.02 & 10.80 & 83.56 & 1.77\\
\sourcefile{PDO9-18W-2A-001} & 2.00 & 18.03 & 20.93 & 86.14 & 2.90\\
\sourcefile{PDO9-27W-3A-001} & 3.00 & 27.00 & 31.10 & 86.82 & 4.10\\
\sourcefile{PDO15-15W-1A-001} & 1.00 & 15.01 & 17.37 & 86.41 & 2.36\\
\bottomrule
\end{longtable}
\endgroup

% Automatically rebuilt from original records by rebuild_results.py.
\begingroup
\scriptsize
\setlength{\tabcolsep}{3pt}
\renewcommand{\arraystretch}{1.12}
\begin{longtable}{@{}L{0.2728\linewidth} L{0.1408\linewidth} L{0.1056\linewidth} L{0.0704\linewidth} L{0.1056\linewidth} L{0.1848\linewidth}@{}}
\caption{Power factor and current THD from the instrument readings. The instruments or averaging intervals can differ. The 3.9 A entry is an associated record and is not used for the load trend.}\label{tab:pq-results}\\
\toprule
Test ID & Setting / load & $P_\mathrm{in}$ (W) & PF & Current THD (\%) & Instrument / record\\
\midrule
\endfirsthead
\toprule
Test ID & Setting / load & $P_\mathrm{in}$ (W) & PF & Current THD (\%) & Instrument / record\\
\midrule
\endhead
\sourcefile{PDO5-5W-001} & 5 V / 1.00 A & 9.195 & -- & 62.10 & ITECH entries\\
\sourcefile{PDO5-10W-2A-001} & 5 V / 2.00 A & 14.779 & 0.620 & 28.37 & ITECH entries\\
\sourcefile{PDO5-MAX3A-001} & 5 V / 3.00 A & 21.585 & 0.620 & 40.80 & ITECH entries\\
\sourcefile{PDO20-5W-001} & 20 V / 0.25 A & 8.780 & 0.440 & 79.46 & ITECH entries\\
\sourcefile{PDO20-75W-001} & 20 V / 3.70 A & 95.875 & 0.970 & 2.23 & ITECH entries\\
\sourcefile{PDO20-76W-38A-001} & 20 V / 3.80 A & 98.990 & 0.980 & 2.62 & ITECH entries\\
\sourcefile{PDO20-80W-001} & 20 V / 3.90 A & 95.875 & 0.970 & 2.23 & Associated record\\
\sourcefile{PDO20-80W-40A-001} & 20 V / 4.00 A & 104.418 & 0.990 & 3.98 & ITECH entries\\
\sourcefile{PDO9-9W-1A-001} & 9 V / 1.00 A & 10.795 & 0.223 & -- & Siglent PF; ITECH THD\\
\sourcefile{PDO9-18W-2A-001} & 9 V / 2.00 A & 20.927 & 0.364 & 4.38 & Siglent PF; ITECH THD\\
\sourcefile{PDO9-27W-3A-001} & 9 V / 3.00 A & 31.101 & 0.474 & 6.81 & Siglent PF; ITECH THD\\
\sourcefile{PDO15-15W-1A-001} & 15 V / 1.00 A & 17.372 & 0.316 & 12.43 & Siglent PF; ITECH THD\\
\sourcefile{PDO15-30W-2A-001} & 15 V / 2.00 A & 34.071 & 0.499 & 45.00 & Siglent PF; ITECH THD\\
\sourcefile{PDO15-45W-3A-001} & 15 V / 3.00 A & 50.622 & 0.638 & 72.20 & Siglent PF; ITECH THD\\
\sourcefile{PDO20-94W-47A-AC-001} & 20 V / 4.70 A & 102.265 & 0.894 & 32.42 & Siglent PF; ITECH THD\\
\bottomrule
\end{longtable}
\endgroup

% Automatically rebuilt from original records by rebuild_results.py.
\begingroup
\scriptsize
\setlength{\tabcolsep}{3pt}
\renewcommand{\arraystretch}{1.12}
\begin{longtable}{@{}L{0.2694\linewidth} L{0.0898\linewidth} L{0.1257\linewidth} L{0.1257\linewidth} L{0.1347\linewidth} L{0.1347\linewidth}@{}}
\caption{Load-removal measurements from the workbook. Averages and pp span the whole capture and do not give the voltage rise from the pre-step level. Recovery time is not available. The 5 V and 20 V screenshots have no numerical worksheet row.}\label{tab:appendix-transient-complete}\\
\toprule
Test ID & Setting (V) & VBUS avg (V) & VBUS pp (V) & Secondary avg (V) & Secondary pp (V)\\
\midrule
\endfirsthead
\toprule
Test ID & Setting (V) & VBUS avg (V) & VBUS pp (V) & Secondary avg (V) & Secondary pp (V)\\
\midrule
\endhead
\sourcefile{PDO9-3A-TO-NL-001} & 9 & 9.08221 & 1.145 & 31.0671 & 25.1276\\
\sourcefile{PDO15-3A-TO-NL-001} & 15 & 15.01610 & 0.414 & 30.9675 & 25.2539\\
\bottomrule
\end{longtable}
\endgroup

% Automatically rebuilt from original records by rebuild_results.py.
\begin{landscape}
\begingroup
\scriptsize
\setlength{\tabcolsep}{3pt}
\renewcommand{\arraystretch}{1.12}
\begin{longtable}{@{}L{0.2604\linewidth} L{0.1347\linewidth} L{0.2155\linewidth} L{0.2694\linewidth}@{}}
\caption{Test status, firmware version, and use in the calculations. A dash means the firmware was not recorded. Efficiency estimates are preliminary because acquisition timing and inclusion of auxiliary power were not fully documented. AC-coupled statistics are not used to calculate DC power.}\label{tab:appendix-record-metadata}\\
\toprule
Test ID & Recorded status & Firmware & Treatment\\
\midrule
\endfirsthead
\toprule
Test ID & Recorded status & Firmware & Treatment\\
\midrule
\endhead
\sourcefile{PDO5-NL-001} & Measured & PWR\_1\_9\_20 & No nonzero DC input/output power pair\\
\sourcefile{PDO5-5W-001} & Measured & PWR\_1\_9\_20 & Power comparison\\
\sourcefile{PDO5-10W-2A-001} & Measured & PWR\_1\_9\_20 & Power comparison\\
\sourcefile{PDO5-MAX3A-001} & Measured & PWR\_1\_9\_20 & Power comparison\\
\sourcefile{PDO9-NL-001} & Measured & -- & No nonzero DC input/output power pair\\
\sourcefile{PDO15-NL-001} & Measured & -- & No nonzero DC input/output power pair\\
\sourcefile{PDO20-5W-001} & Measured & PWR\_1\_9\_20 & Power comparison\\
\sourcefile{PDO20-75W-001} & Measured & PWR\_1\_9\_20 & Power comparison\\
\sourcefile{PDO20-76W-38A-001} & Measured & PWR\_1\_9\_20 & Power comparison\\
\sourcefile{PDO20-80W-001} & Measured & PWR\_1\_9\_20 & associated input-power record\\
\sourcefile{PDO20-80W-40A-001} & Measured & PWR\_1\_9\_20 & Power comparison\\
\sourcefile{PDO9-9W-1A-001} & Measured & -- & Power comparison\\
\sourcefile{PDO9-18W-2A-001} & Measured & -- & Power comparison\\
\sourcefile{PDO9-27W-3A-001} & Measured & -- & Power comparison\\
\sourcefile{PDO9-3A-TO-NL-001} & Captured & -- & transient\\
\sourcefile{PDO15-15W-1A-001} & Measured & -- & Power comparison\\
\sourcefile{PDO15-30W-2A-001} & Captured & -- & instrument power disagreement\\
\sourcefile{PDO15-45W-3A-001} & Recorded & -- & instrument power disagreement\\
\sourcefile{PDO15-3A-TO-NL-001} & Recorded & -- & transient\\
\sourcefile{PDO15-NL-AC-001} & Recorded & -- & AC-coupled: no DC output\\
\sourcefile{PDO15-15W-1A-AC-001} & Recorded & -- & AC-coupled: no DC output\\
\sourcefile{PDO9-18W-2A-THERM-001} & Thermal screenshot only & -- & No nonzero DC input/output power pair\\
\sourcefile{PDO15-30W-2A-AC-001} & Recorded & -- & AC-coupled: no DC output\\
\sourcefile{PDO15-45W-3A-AC-001} & Recorded & -- & AC-coupled: no DC output\\
\sourcefile{PDO5-5W-1A-AC-001} & Recorded & -- & AC-coupled: no DC output\\
\sourcefile{PDO5-10W-2A-AC-001} & Recorded & -- & AC-coupled: no DC output\\
\sourcefile{PDO5-15W-3A-AC-001} & Recorded & -- & AC-coupled: no DC output\\
\sourcefile{PDO20-NL-AC-001} & Recorded & -- & AC-coupled: no DC output\\
\sourcefile{PDO20-40W-2A-AC-001} & Recorded & -- & AC-coupled: no DC output\\
\sourcefile{PDO20-60W-3A-AC-001} & Recorded & -- & AC-coupled: no DC output\\
\sourcefile{PDO20-80W-4A-AC-001} & Recorded & -- & AC-coupled: no DC output\\
\sourcefile{PDO20-94W-47A-AC-001} & Recorded & PWR-1.9.10 DMA bang-bang controller & AC-coupled: no DC output\\
\bottomrule
\end{longtable}
\endgroup
\end{landscape}



\section{Current Harmonic Spectra}
% Automatically rebuilt from original records by rebuild_results.py.
\begingroup
\scriptsize
\setlength{\tabcolsep}{3pt}
\renewcommand{\arraystretch}{1.12}
\begin{longtable}{@{}L{0.1135\linewidth} L{0.2555\linewidth} L{0.2555\linewidth} L{0.2555\linewidth}@{}}
\caption{Input-current harmonics from the three ITECH tables. Order 1 is the fundamental. The 4.7 A entry uses an AC-coupled output capture.}\label{tab:input-harmonics}\\
\toprule
Order & 3.8 A: \% of fundamental & 4.0 A: \% of fundamental & 4.7 A: \% of fundamental\\
\midrule
\endfirsthead
\toprule
Order & 3.8 A: \% of fundamental & 4.0 A: \% of fundamental & 4.7 A: \% of fundamental\\
\midrule
\endhead
1 & 100.0 & 100.0 & 100.0\\
2 & 0.2 & 2.6 & 0.1\\
3 & 0.8 & 2.3 & 12.9\\
4 & 0.2 & 1.0 & 0.0\\
5 & 0.7 & 1.2 & 16.3\\
6 & 0.1 & 0.5 & 0.1\\
7 & 0.6 & 0.8 & 14.7\\
8 & 0.1 & 0.4 & 0.1\\
9 & 0.1 & 0.3 & 12.3\\
10 & 0.1 & 0.2 & 0.1\\
11 & 0.4 & 0.5 & 9.3\\
12 & 0.1 & 0.2 & 0.1\\
13 & 0.1 & 0.2 & 7.2\\
14 & 0.1 & 0.1 & 0.0\\
15 & 0.4 & 0.3 & 5.1\\
16 & 0.1 & 0.2 & 0.0\\
17 & 0.1 & 0.2 & 3.8\\
18 & 0.2 & 0.0 & 0.0\\
19 & 0.3 & 0.2 & 3.6\\
20 & 0.1 & 0.3 & 0.0\\
21 & 0.5 & 0.3 & 3.6\\
22 & 1.5 & 1.3 & 0.0\\
23 & 0.8 & 0.5 & 3.5\\
24 & 1.2 & 1.3 & 0.0\\
25 & 0.3 & 0.4 & 3.5\\
26 & 0.2 & 0.3 & 0.0\\
27 & 0.2 & 0.2 & 3.1\\
28 & 0.1 & 0.1 & 0.0\\
29 & 0.1 & 0.1 & 2.6\\
30 & 0.1 & 0.1 & 0.0\\
31 & 0.1 & 0.2 & 2.2\\
32 & 0.1 & 0.1 & 0.0\\
33 & 0.0 & 0.2 & 1.6\\
34 & 0.1 & 0.1 & 0.0\\
35 & 0.1 & 0.1 & 1.2\\
36 & 0.0 & 0.1 & 0.1\\
37 & 0.1 & 0.1 & 1.0\\
38 & 0.0 & 0.0 & 0.1\\
39 & 0.0 & 0.1 & 0.8\\
40 & 0.1 & 0.0 & 0.1\\
\bottomrule
\end{longtable}
\endgroup



\section{Thermal-Camera Readings}
% Automatically rebuilt from original records by rebuild_results.py.
\begingroup
\scriptsize
\setlength{\tabcolsep}{3pt}
\renewcommand{\arraystretch}{1.12}
\begin{longtable}{@{}L{0.3400\linewidth} L{0.1800\linewidth} L{0.1800\linewidth} L{0.1800\linewidth}@{}}
\caption{Temperature readings from ten thermal images, using emissivity 0.95. The lower reference marker was not confirmed as ambient temperature.}\label{tab:thermal-results-complete}\\
\toprule
Image / condition & Maximum (°C) & Cursor (°C) & Reference (°C)\\
\midrule
\endfirsthead
\toprule
Image / condition & Maximum (°C) & Cursor (°C) & Reference (°C)\\
\midrule
\endhead
Image \#1 / 20 V, 4.7 A & 102.9 & 27.9 & 25.2\\
Image \#2 / 20 V, 1 A & 60.3 & 30.8 & 24.9\\
Image \#3 / 20 V, 3 A & 82.9 & 30.6 & 25.2\\
Image \#4 / 5 V, 1 A & 47.8 & 29.1 & 24.7\\
Image \#5 / 20 V, no load & 38.0 & 29.2 & 25.4\\
Image \#6 / 20 V, 4 A & 90.4 & 31.5 & 24.5\\
Image \#7 / 5 V, open & 37.5 & 34.9 & 25.5\\
Image \#8 / 9 V, 1 A & 49.9 & 29.6 & 25.0\\
Image \#9 / 9 V, 2 A & 56.8 & 39.7 & 24.9\\
Image \#10 / 9 V, 3 A & 78.0 & 32.1 & 25.3\\
\bottomrule
\end{longtable}
\endgroup



\section{Waveform Acquisition}
% Automatically rebuilt from original records by rebuild_results.py.
\begingroup
\scriptsize
\setlength{\tabcolsep}{3pt}
\renewcommand{\arraystretch}{1.12}
\begin{longtable}{@{}L{0.3562\linewidth} L{0.2619\linewidth} L{0.2619\linewidth}@{}}
\caption{Waveform records, sample counts, and sampling intervals.}\label{tab:raw-waveform-sources}\\
\toprule
Record & Samples & Interval (s)\\
\midrule
\endfirsthead
\toprule
Record & Samples & Interval (s)\\
\midrule
\endhead
\sourcefile{00006} & 500002 & 4e-09\\
\sourcefile{00007} & 12502 & 4e-10\\
\sourcefile{00008} & 500002 & 2e-09\\
\sourcefile{00010} & 50002 & 4e-10\\
\sourcefile{lecory 7 C1} & 10001 & 1e-05\\
\sourcefile{lecory 7 C2} & 10001 & 1e-05\\
\sourcefile{lecory 7 C3} & 10001 & 1e-05\\
\sourcefile{lecory 8 C1} & 10001 & 1e-05\\
\sourcefile{lecory 8 C2} & 10001 & 1e-05\\
\sourcefile{lecory 8 C3} & 10001 & 1e-05\\
\bottomrule
\end{longtable}
\endgroup

